\chapter{Pending Generation Experiments}\label{ch:stats}
No inferential statistical tests have been performed for generated captions because no generation outputs exist. The project reports descriptive counts, means, distributions, and a retrieval similarity mean, but it has no paired metric vectors, human ratings, or method-level outcomes from which to calculate confidence intervals or significance tests. Accordingly, no p-values or claims of statistical significance are reported.

For the future paired comparison, each held-out image should be evaluated under every RAG condition with identical references. Report per-image metric differences and bootstrap confidence intervals. For metric distributions that support paired parametric inference, a paired test may be reported with effect size; otherwise use a paired nonparametric test such as Wilcoxon signed-rank, with correction for multiple comparisons. Human judgments should be analyzed with an appropriate repeated-measures or mixed-effects model. These are recommendations, not completed analyses.

\section{Ablation Analysis}\label{ch:ablation}
An ablation requires measured outputs under controlled variants. No caption generator or RAG method is implemented, so changing retrieval depth, removing retrieved evidence, altering the retrieved modality, or varying evidence formatting cannot yield valid caption-quality results in this repository. The geometry-only preprocessing comparison in \cref{tab:preprocess} is not a model ablation: it evaluates image geometry and does not measure a downstream task.

\begin{table}[htbp]\centering\caption{Planned ablation matrix; no metric values are available.}\label{tab:ablation}
\begin{tabularx}{\textwidth}{lXXX}\toprule Variant & Controlled change & Evaluation & Status\\\midrule No retrieval & Generator only & Caption metrics and grounding & \pending\\Image neighbors & Image-vector top-$k$ evidence & Same held-out set & \pending\\Caption neighbors & Text-vector top-$k$ evidence & Same held-out set & \pending\\Hybrid evidence & Image and caption evidence & Same held-out set & \pending\\Top-$k$ sensitivity & Vary $k$, fixed generator & Quality/faithfulness trade-off & \pending\\\bottomrule\end{tabularx}\end{table}

\section{Hallucination Analysis}\label{ch:hallucination}
Hallucination analysis should test whether a generated caption asserts objects, attributes, counts, actions, or relations unsupported by the image. The current project has no generated captions and no hallucination annotations, so a hallucination rate cannot be calculated. COCO category labels and boxes may support object-presence checks but do not provide exhaustive scene truth: an unannotated object is not necessarily absent, and a box does not establish an action or relation. CLIP similarity alone is not a reliable factuality validator.

A future protocol should define claim categories before evaluation, annotate a stratified sample with independent raters, assess agreement, and compare generated claims against image evidence. Automated object consistency can be reported separately from human judgments and should include known annotation-coverage limitations. No fabricated hallucination scores are included.

\section{Qualitative Analysis}\label{ch:qualitative}
The repository includes qualitative retrieval records and visual audit materials, but it does not produce captions for qualitative comparison. Retrieval neighbors can be inspected to identify plausible visual analogies and failure cases, yet similarity scores do not prove that the retrieved caption is appropriate for the query. No fabricated image-caption examples are presented as experimental outputs.

Future qualitative review should sample successes and failures across scene complexity and object-frequency strata. It should show the query image, retrieved images/captions, generated output, reference captions, and a concise judgment of evidence support. The existing analyses indicate useful dimensions for this review: object count, category mix, color family, scene predictions, and caption lexical patterns. Zero-shot scene tags and inferred geometric relations should be visibly marked as model-derived or computed metadata.
